\section{Results}
\label{sec:results}

\subsection{Current status: no verified benchmark results}

At the time of writing, no condition in Table~\ref{tab:conditions} has completed a
research-scale training and evaluation cycle. Consequently, this section reports
\textbf{no performance numbers}. The tables below are reporting templates. Dashes
and \pending{} markers are intentional. They must not be cited as
findings, and they must not be replaced with borrowed numbers from other models
or blog posts.

\begin{table}[t]
\centering
\small
\caption{Primary results template. Every cell is pending. A completed table must
include dataset versions, prompt templates, seeds, hardware, and raw
predictions.}
\label{tab:results-main}
\setlength{\tabcolsep}{4pt}
\begin{tabular}{@{}lrrrrr@{}}
\toprule
Metric & Base & LoRA & QLoRA & Distilled & +RAG \\
\midrule
Held-out security MCQ accuracy & \pending{} & \pending{} & \pending{} & \pending{} & \pending{} \\
MCQ 95\% CI & \pending{} & \pending{} & \pending{} & \pending{} & \pending{} \\
Vulnerability macro F1 & \pending{} & \pending{} & \pending{} & \pending{} & \pending{} \\
Vulnerability FPR & \pending{} & \pending{} & \pending{} & \pending{} & \pending{} \\
Secure-code exact match & \pending{} & \pending{} & \pending{} & \pending{} & \pending{} \\
Secure-code syntax valid & \pending{} & \pending{} & \pending{} & \pending{} & \pending{} \\
Grounded QA token F1 & \pending{} & \pending{} & \pending{} & \pending{} & \pending{} \\
Grounding proxy (lower better) & \pending{} & \pending{} & \pending{} & \pending{} & \pending{} \\
General retention accuracy & \pending{} & \pending{} & \pending{} & \pending{} & \pending{} \\
\bottomrule
\end{tabular}
\end{table}

\begin{table}[t]
\centering
\small
\caption{Operational results template. Peak memory and latency are measured on
the stated hardware, not inferred from model size.}
\label{tab:results-ops}
\begin{tabular}{@{}lrrrrr@{}}
\toprule
Measurement & Base & LoRA & QLoRA & Distilled & 4-bit \\
\midrule
Peak training VRAM (GB) & \pending{} & \pending{} & \pending{} & \pending{} & \pending{} \\
Training duration (hours) & \pending{} & \pending{} & \pending{} & \pending{} & \pending{} \\
Peak inference memory (GB) & \pending{} & \pending{} & \pending{} & \pending{} & \pending{} \\
Tokens/second (CPU) & \pending{} & \pending{} & \pending{} & \pending{} & \pending{} \\
Tokens/second (GPU) & \pending{} & \pending{} & \pending{} & \pending{} & \pending{} \\
Cost at stated rate (USD) & \pending{} & \pending{} & \pending{} & \pending{} & \pending{} \\
\bottomrule
\end{tabular}
\end{table}

\subsection{Implementation validation (not a scientific result)}

The following engineering facts were verified during repository construction on a
6-thread Intel i5-9500T CPU with 15.4 GiB RAM, no GPU, PyTorch 2.15.0+CPU,
Transformers 5.24.1, PEFT 0.22.0, and Python 3.12:

\begin{itemize}[leftmargin=1.4em]
  \item the checkpoint loaded as \texttt{Qwen3\_5ForConditionalGeneration} with
  852,985,920 parameters;
  \item a rank-4 LoRA adapter added 2.71 million trainable parameters after
  freezing the vision tower;
  \item a forward/backward pass on a two-example batch produced a finite loss and
  gradients;
  \item a two-step smoke run completed in 26--29 seconds of training after model
  load, used about 6.25 GiB peak RSS, saved an adapter, and reloaded it for
  generation;
  \item the FIM and prefix autocomplete code paths executed without a hidden
  execution path.
\end{itemize}

These checks establish that the plumbing runs. They say nothing about
cybersecurity capability, and the smoke adapter is intentionally undertrained.
A poor chat answer or a poor FIM completion from this adapter is expected; it is
not a result about a trained \cyberlite{} model.

\subsection{What will count as a result}

A result will be promoted only after:

\begin{enumerate}[leftmargin=1.4em]
  \item the training run has a clean commit, config hash, and dataset manifest;
  \item the evaluation uses a fixed, licensed held-out set that was not used for
  tuning;
  \item raw predictions are archived with the summary;
  \item the hardware, software, decoding, and seeds are recorded;
  \item the authors have reviewed a sample for unsafe or unsupported claims;
  \item the table includes negative or null findings, not only improvements.
\end{enumerate}

\subsection{Expected failure modes}

We anticipate several possible outcomes, all of which are reportable:
the 0.8B model may improve on knowledge questions but remain weak at multi-step
reasoning; QLoRA may match LoRA on some tasks but diverge on others; retrieval
may improve grounding while adding latency; FIM may fail if the base checkpoint
was not FIM-trained; and narrow SFT may degrade general capability. These are
hypotheses, not predictions disguised as results.
